\clearpage
\begingroup
\fancyhead[L]{}
\renewcommand{\headrulewidth}{0pt}
\SourceLeaf{646}
\needspace{6\baselineskip}
\pdfbookmark[1]{Remarks on the Theory of the Motion of the Moon}{moon-bemerkungen}
\begin{center}\large REMARKS ON THE THEORY OF THE MOTION OF THE MOON.\end{center}

\par\noindent\hspace*{\parindent}The »Theorie der Bewegung des Mondes« printed above, which is written on a bundle of
sheets, remained unfinished. Besides the manuscript printed here, the posthumous papers contain two
earlier, less developed drafts. In \textsc{Gauss}’s scientific diary, there is the following note under 1801
August: »\textit{Theoriam motus Lunae aggressi sumus}«. The manuscript must in fact date from this period.
In addition to the handwriting and the manner of presentation, this early dating is supported by the
fact that the anomalies are reckoned from the apogee rather than the perigee. Compare
the perturbations of Ceres in this connection.

\par\noindent\hspace*{\parindent}The interest of these investigations by \textsc{Gauss}, in comparison with the methods of other authors, including
later ones, lies in his distinctive treatment of the differential equations of the Moon’s motion.
The integration and the determination of the perturbation terms remained incomplete.
Only the latitude perturbations are complete, and \textsc{Gauss} evidently never returned to these papers.
He did not comment on them anywhere else, even in the correspondence beginning in 1802
with \textsc{Olbers}. As early as 1802, the third volume of \textsc{Laplace’s}
Mécanique céleste appeared, containing the lunar theory and superseding the developments \textsc{Gauss}
had begun. This circumstance, as well as the direct connection with the work of \textsc{Tobias Mayer}, supports the
early date mentioned above.

\par\noindent\hspace*{\parindent}The form in which \textsc{Gauss} ultimately presents the perturbations is the same as \textsc{Plana}’s
(1832), since both expand the divisors in powers of the quantity \(n\) (the ratio of the orbital periods of the Moon and
the Sun).

\par\noindent\hspace*{\parindent}Some small additions and corrections have been made to present the results completely and correctly, at least
within the limits initially adopted by \textsc{Gauss}. Thus, in integrating
the second equation of the third section, p.~622, the terms in \(\varepsilon\sin(2E\pm a)\) were added and
retained in the following equations, although, as \textsc{Gauss} notes, they can be omitted there
because even after the integrations they are still only of the third order. Their inclusion is, however,
necessary for the second approximation to obtain the correct coefficients of the analogous terms in the expression
for \(\frac{\mathrm du^2}{\mathrm dv^2}\pi\), p.~634.

\par\noindent\hspace*{\parindent}In the fourth section the differential equation for \(\theta\) is accurate up to and including the terms of order~5.
After integration, the terms that the divisor \(n\) brings into
order~5 are therefore naturally missing, as is the term in \(\sin(2E-\lambda+2a)\) that receives the divisor \(nn\).
\textsc{Gauss} marks its omission with an \(*\). According to \textsc{Plana}, the term equals \(-\frac{9}{32}n\varepsilon\varepsilon\sin(2E-\lambda+2a)\).
In the transformation to \textsc{Mayer}’s form, pp.~631---633, the lists of terms to be added,
which are themselves incomplete, were omitted in printing. The comparison of the results of
\textsc{Gauss} and \textsc{Mayer} was also left without a check calculation.

\VIIInlineSourceHead{\SourceLeaf{647}}{}{REMARKS.\quad THEORY OF THE MOTION OF THE MOON.}{641}
\par\noindent\hspace*{\parindent}In the fifth section \textsc{Gauss} set up the differential equation for \(p\) only as far as the terms of the fourth order
inclusive. In expanding the expression for \(\sin 2(v-V)\), however, some
terms of the third order were retained because they acquire small divisors on integration. The two
terms in brackets are added here solely because \textsc{Gauss} took them into account in setting up the
following expression. Naturally, the final expression for \(p\) omits
the terms below order~5 that arise on integrating terms of order~5.
\textsc{Gauss} had, it is true, later added some of these terms to his manuscript. But
since he had done so only sporadically and hastily, and in many cases the values given were still incorrect,
these terms were omitted from the printed text. Some other inaccuracies were also found
and corrected. \textsc{Gauss} evidently broke off his investigations here rather abruptly,
without having checked his calculations.
\begin{flushright}\textsc{Brendel}.\end{flushright}
\par\begin{center}\rule{.24\linewidth}{.4pt}\end{center}
\clearpage
\endgroup
